% Options for packages loaded elsewhere
\PassOptionsToPackage{unicode}{hyperref}
\PassOptionsToPackage{hyphens}{url}
\PassOptionsToPackage{dvipsnames,svgnames,x11names}{xcolor}
\documentclass[
  10pt,
  a4paper,
]{article}
\usepackage{xcolor}
\usepackage[margin=1.8cm]{geometry}
\usepackage{amsmath,amssymb}
\setcounter{secnumdepth}{5}
\usepackage{iftex}
\ifPDFTeX
  \usepackage[T1]{fontenc}
  \usepackage[utf8]{inputenc}
  \usepackage{textcomp} % provide euro and other symbols
\else % if luatex or xetex
  \usepackage{unicode-math} % this also loads fontspec
  \defaultfontfeatures{Scale=MatchLowercase}
  \defaultfontfeatures[\rmfamily]{Ligatures=TeX,Scale=1}
\fi
\usepackage{lmodern}
\ifPDFTeX\else
  % xetex/luatex font selection
  \setmainfont[]{Arial}
  \setsansfont[]{Arial}
  \setmonofont[]{/Users/teobale/Library/Fonts/JetBrainsMonoNL-Regular.ttf}
\fi
% Use upquote if available, for straight quotes in verbatim environments
\IfFileExists{upquote.sty}{\usepackage{upquote}}{}
\IfFileExists{microtype.sty}{% use microtype if available
  \usepackage[]{microtype}
  \UseMicrotypeSet[protrusion]{basicmath} % disable protrusion for tt fonts
}{}
\usepackage{setspace}
\makeatletter
\@ifundefined{KOMAClassName}{% if non-KOMA class
  \IfFileExists{parskip.sty}{%
    \usepackage{parskip}
  }{% else
    \setlength{\parindent}{0pt}
    \setlength{\parskip}{6pt plus 2pt minus 1pt}}
}{% if KOMA class
  \KOMAoptions{parskip=half}}
\makeatother
\usepackage{longtable,booktabs,array}
\newcounter{none} % for unnumbered tables
\usepackage{calc} % for calculating minipage widths
% Correct order of tables after \paragraph or \subparagraph
\usepackage{etoolbox}
\makeatletter
\patchcmd\longtable{\par}{\if@noskipsec\mbox{}\fi\par}{}{}
\makeatother
% Allow footnotes in longtable head/foot
\IfFileExists{footnotehyper.sty}{\usepackage{footnotehyper}}{\usepackage{footnote}}
\makesavenoteenv{longtable}
\setlength{\emergencystretch}{3em} % prevent overfull lines
\providecommand{\tightlist}{%
  \setlength{\itemsep}{0pt}\setlength{\parskip}{0pt}}
\setlength{\emergencystretch}{3em}
\widowpenalty=10000
\clubpenalty=10000
\makeatletter
\g@addto@macro\verbatim\small
\makeatother
\usepackage{bookmark}
\IfFileExists{xurl.sty}{\usepackage{xurl}}{} % add URL line breaks if available
\urlstyle{same}
% fallback for those not using the hyperref driver hyperxmp:
\makeatletter
\@ifundefined{xmpquote}{\newcommand{\xmpquote}[1]{#1}}{}
\makeatother
\hypersetup{
  pdftitle={Basi di dati - Appunti mirati all'esame},
  colorlinks=true,
  linkcolor={blue},
  filecolor={Maroon},
  citecolor={Blue},
  urlcolor={blue},
  pdfcreator={LaTeX via pandoc}}

\title{Basi di dati - Appunti mirati all'esame}
\usepackage{etoolbox}
\makeatletter
\providecommand{\subtitle}[1]{% add subtitle to \maketitle
  \apptocmd{\@title}{\par {\large #1 \par}}{}{}
}
\makeatother
\subtitle{Sintesi exam-first basata sugli appelli 2024-2025}
\author{}
\date{Giugno 2026}

\begin{document}
\maketitle

{
\setcounter{tocdepth}{2}
\tableofcontents
}
\setstretch{1.06}
Questi appunti sono costruiti sui cinque appelli disponibili
(16/07/2024, 14/01/2025, 10/06/2025, 01/07/2025, 22/07/2025). Il
programma del corso è usato per chiarire e completare ciò che compare
nelle prove, non per assegnare la stessa importanza a ogni argomento.

Per gli svolgimenti completi si veda \url{SOLUZIONI_APPELLI.md}.

\section{1. Cosa studiare prima}\label{1-cosa-studiare-prima}

{\def\LTcaptype{none} % do not increment counter
\begin{longtable}[]{@{}lrr@{}}
\toprule\noalign{}
Argomento & Presenza nei 5 appelli & Priorità \\
\midrule\noalign{}
\endhead
\bottomrule\noalign{}
\endlastfoot
SQL (join, negazione, aggregazione, massimo, divisione) & 5/5 &
massima \\
Algebra relazionale & 5/5 & massima \\
ER, traduzione e reverse engineering & 5/5 & massima \\
Vincoli SQL, chiavi, \texttt{NULL}, effetti delle modifiche & 5/5 &
massima \\
Normalizzazione & 5/5 & alta \\
Sicurezza, \texttt{GRANT}/\texttt{REVOKE} & 5/5 & alta \\
Transazioni e proprietà ACID & 3/5 & media-alta \\
File e indici & 3/5 & media-alta \\
Trigger, viste, PL/pgSQL, NoSQL/MongoDB & 0/5 & bassa, ripasso finale \\
\end{longtable}
}

Negli appelli estivi 2025 l\textquotesingle esercizio 1 vale 8 punti ed
è una \textbf{soglia}: se non si raggiungono tutti gli 8 punti, il resto
può non essere corretto. Va quindi preparato come una batteria di
risposte brevi e molto precise.

\section{2. Strategia per i 90
minuti}\label{2-strategia-per-i-90-minuti}

\begin{enumerate}
\def\labelenumi{\arabic{enumi}.}
\tightlist
\item
  Leggere DDL, chiavi, \texttt{NOT\ NULL}, azioni
  \texttt{ON\ DELETE/UPDATE} e annotare le FK.
\item
  Svolgere con calma l\textquotesingle esercizio-soglia; controllare
  soprattutto \texttt{NULL}, join esterni, duplicati e vincoli impliciti
  della chiave primaria.
\item
  Fare prima le query SQL riconducibili a un pattern noto.
\item
  Scrivere l\textquotesingle algebra con relazioni intermedie e
  ridenominazioni esplicite.
\item
  Disegnare ER/logico e teoria breve con definizioni, non con esempi
  vaghi.
\item
  Lasciare almeno 8-10 minuti per verificare alias, \texttt{GROUP\ BY},
  casi a zero e direzione delle chiavi esterne.
\end{enumerate}

\begin{center}\rule{0.5\linewidth}{0.5pt}\end{center}

Parte I - Modello relazionale, DDL e domande rapide

\section{3. Relazione, schema e
istanza}\label{3-relazione-schema-e-istanza}

\begin{itemize}
\tightlist
\item
  Una relazione è un insieme di tuple su domini atomici.
\item
  \textbf{Grado}: numero di attributi.
\item
  \textbf{Cardinalità}: numero di tuple dell\textquotesingle istanza.
\item
  L\textquotesingle ordine di righe e colonne non ha significato nel
  modello relazionale.
\item
  Nel modello relazionale puro non esistono tuple duplicate; SQL,
  invece, usa normalmente semantica a multinsieme e può produrre
  duplicati.
\end{itemize}

\section{4. Superchiave, chiave e chiave
primaria}\label{4-superchiave-chiave-e-chiave-primaria}

\begin{itemize}
\tightlist
\item
  \textbf{Superchiave}: insieme di attributi che identifica univocamente
  ogni tupla.
\item
  \textbf{Chiave candidata}: superchiave minimale. Togliendo un
  attributo perde l\textquotesingle unicità.
\item
  \textbf{Chiave primaria}: una chiave candidata scelta come
  identificatore principale. È unica per tabella e implica
  \texttt{UNIQUE} e \texttt{NOT\ NULL}.
\item
  Una superchiave non minimale si ottiene aggiungendo attributi a una
  chiave: se \texttt{code} è PK, \texttt{\{code,\ air\_company\}} è una
  superchiave ma non una chiave.
\item
  \textbf{Attributo primo}: attributo che appartiene ad almeno una
  chiave candidata.
\end{itemize}

Attenzione: \texttt{UNIQUE} non equivale sempre a \texttt{PRIMARY\ KEY},
perché in SQL può ammettere valori nulli; la PK non li ammette.

\section{5. Chiavi esterne e integrità
referenziale}\label{5-chiavi-esterne-e-integrituxe0-referenziale}

Una FK \texttt{R.X\ REFERENCES\ S.Y} impone che ogni valore non nullo di
\texttt{R.X} compaia nel campo candidato referenziato \texttt{S.Y}.

\begin{itemize}
\tightlist
\item
  La FK è nullable se non ha \texttt{NOT\ NULL} e non appartiene a una
  PK.
\item
  Gli attributi di una PK composta sono tutti implicitamente
  \texttt{NOT\ NULL}.
\item
  Se non è specificata un\textquotesingle azione,
  \texttt{ON\ DELETE/UPDATE} è normalmente \texttt{NO\ ACTION}:
  l\textquotesingle operazione che spezzerebbe i riferimenti viene
  rifiutata.
\item
  \texttt{CASCADE}: propaga eliminazione o aggiornamento alle tuple
  referenzianti.
\item
  \texttt{SET\ NULL}: mette a \texttt{NULL} la FK; è possibile solo se
  la colonna è nullable.
\end{itemize}

Un\textquotesingle istruzione è atomica. Se un \texttt{DELETE} attiva un
\texttt{CASCADE} su una tabella ma viola una FK \texttt{NO\ ACTION} su
un\textquotesingle altra, fallisce \textbf{l\textquotesingle intera
istruzione}: non resta neppure la cancellazione in cascata parziale.

\section{\texorpdfstring{6. \texttt{NULL}, \texttt{CHECK} e lettura del
DDL}{6. NULL, CHECK e lettura del DDL}}\label{6-null-check-e-lettura-del-ddl}

\begin{itemize}
\tightlist
\item
  \texttt{NULL} significa valore mancante; \texttt{NULL\ =\ NULL} non è
  vero ma \texttt{UNKNOWN}.
\item
  Si testa con \texttt{IS\ NULL} / \texttt{IS\ NOT\ NULL}, mai con
  \texttt{=\ NULL}.
\item
  Una condizione \texttt{WHERE} conserva solo le righe per cui il
  predicato è \texttt{TRUE}.
\item
  Un \texttt{CHECK} è violato solo quando è \texttt{FALSE};
  \texttt{TRUE} e \texttt{UNKNOWN} passano. Per vietare il nullo serve
  anche \texttt{NOT\ NULL}.
\item
  Una stringa vuota \texttt{\textquotesingle{}\textquotesingle{}} non è
  \texttt{NULL} in PostgreSQL.
\item
  Non inventare vincoli semantici non presenti nel DDL. Per esempio,
  \texttt{expiry\_date\ \textless{}\ release\_date} è sospetto, ma non
  viola il DDL se manca il relativo \texttt{CHECK}.
\end{itemize}

Checklist rapida su un\textquotesingle istanza:

\begin{enumerate}
\def\labelenumi{\arabic{enumi}.}
\tightlist
\item
  tipo e dominio;
\item
  \texttt{NOT\ NULL} e PK;
\item
  unicità di PK/\texttt{UNIQUE};
\item
  \texttt{CHECK};
\item
  FK, ma solo se è visibile anche la tabella referenziata.
\end{enumerate}

\section{7. Join esterni: il punto più facile da
sbagliare}\label{7-join-esterni-il-punto-piuxf9-facile-da-sbagliare}

\begin{verbatim}
FROM R LEFT JOIN S
  ON condizione_di_join AND filtro_su_S
\end{verbatim}

Il filtro nell\textquotesingle{}\texttt{ON} decide quali righe di
\texttt{S} si abbinano; non elimina le righe di \texttt{R}, che sono
preservate dal \texttt{LEFT\ JOIN}.

\begin{verbatim}
FROM R LEFT JOIN S ON ...
WHERE S.id IS NULL
\end{verbatim}

è un \textbf{anti-join}: conserva le righe di \texttt{R} senza
corrispondenza.

Se una riga di \texttt{R} trova due righe di \texttt{S},
\texttt{SELECT\ R.*} restituisce due copie identiche sotto la semantica
SQL, a meno di usare \texttt{DISTINCT}.

\begin{center}\rule{0.5\linewidth}{0.5pt}\end{center}

Parte II - SQL: i pattern che coprono quasi tutte le query

\section{8. Regola generale}\label{8-regola-generale}

Prima di scrivere SQL, determinare:

\begin{itemize}
\tightlist
\item
  una riga del risultato rappresenta che cosa?
\item
  quali tabelle servono per filtrare e quali per mostrare attributi?
\item
  vanno conservati anche gli oggetti con conteggio zero?
\item
  si cercano uno, tutti i massimi, oppure almeno un caso?
\item
  il risultato può contenere duplicati?
\end{itemize}

Negli esempi seguenti \texttt{R} è l\textquotesingle entità principale e
\texttt{S} contiene associazioni.

\section{9. ``Non esiste'' / ``mai''}\label{9-non-esiste--mai}

Pattern consigliato:

\begin{verbatim}
SELECT r.*
FROM R AS r
WHERE NOT EXISTS (
    SELECT 1
    FROM S AS s
    WHERE s.rid = r.id
      AND <condizione vietata>
);
\end{verbatim}

Include anche gli oggetti senza alcuna riga in \texttt{S}. È il pattern
di:

\begin{itemize}
\tightlist
\item
  clienti che non hanno ordini in spedizione;
\item
  clienti che non hanno mai preso libri romantici;
\item
  libri senza posizione;
\item
  persone senza documenti validi.
\end{itemize}

Alternativa anti-join:

\begin{verbatim}
SELECT r.*
FROM R AS r
LEFT JOIN S AS s
  ON s.rid = r.id AND <condizione vietata>
WHERE s.rid IS NULL;
\end{verbatim}

Evitare \texttt{NOT\ IN} se la sottoquery può contenere \texttt{NULL}:
un solo nullo può rendere il confronto \texttt{UNKNOWN} per tutte le
righe.

\section{10. Conteggio includendo lo
zero}\label{10-conteggio-includendo-lo-zero}

Il filtro sulla tabella destra deve stare
nell\textquotesingle{}\texttt{ON}:

\begin{verbatim}
SELECT r.id, COUNT(s.id) AS n
FROM R AS r
LEFT JOIN S AS s
  ON s.rid = r.id
 AND <condizione su S>
GROUP BY r.id;
\end{verbatim}

Usare \texttt{COUNT(s.id)}, non \texttt{COUNT(*)}: la riga esterna
prodotta dal \texttt{LEFT\ JOIN} farebbe risultare 1 anche quando
\texttt{S} manca.

Per una somma:

\begin{verbatim}
COALESCE(SUM(d.quantity * p.price), 0) AS totale
\end{verbatim}

\texttt{COALESCE} trasforma la somma nulla degli oggetti senza dettagli
in zero.

\section{11. Massimo con tutti gli ex
aequo}\label{11-massimo-con-tutti-gli-ex-aequo}

\begin{verbatim}
WITH conteggi AS (
    SELECT r.id, COUNT(s.id) AS n
    FROM R AS r
    LEFT JOIN S AS s ON s.rid = r.id
    GROUP BY r.id
)
SELECT *
FROM conteggi
WHERE n = (SELECT MAX(n) FROM conteggi);
\end{verbatim}

Non usare \texttt{ORDER\ BY\ ...\ LIMIT\ 1} se la traccia dice ``la
persona (o persone)'' o se gli ex aequo devono essere restituiti.

\section{12. Più della media}\label{12-piuxf9-della-media}

Prima costruire una riga di conteggio per ogni oggetto, poi calcolare la
media:

\begin{verbatim}
WITH conteggi AS (
    SELECT r.id, COUNT(s.id) AS n
    FROM R AS r
    LEFT JOIN S AS s
      ON s.rid = r.id
     AND <intervallo temporale>
    GROUP BY r.id
)
SELECT *
FROM conteggi
WHERE n > (SELECT AVG(n) FROM conteggi);
\end{verbatim}

Specificare se la media è su tutti gli oggetti o solo su quelli
``attivi''. La prima interpretazione richiede il \texttt{LEFT\ JOIN}; la
seconda può raggruppare solo \texttt{S}.

\section{13. Coppie e ``due righe diverse'': self
join}\label{13-coppie-e-due-righe-diverse-self-join}

\begin{verbatim}
SELECT DISTINCT x1.id, x2.id
FROM X AS x1
JOIN X AS x2
  ON <proprietà comune>
 AND x1.id < x2.id;
\end{verbatim}

\texttt{x1.id\ \textless{}\ x2.id} elimina sia la coppia \texttt{(x,x)}
sia il duplicato simmetrico \texttt{(y,x)}. Se le due righe hanno ruoli
temporali diversi, è più naturale usare
\texttt{x2.year\ =\ x1.year\ +\ 1}.

Pattern d\textquotesingle esame:

\begin{itemize}
\tightlist
\item
  passeggeri sullo stesso volo;
\item
  stesso prodotto e stessa quantità in due ordini distinti;
\item
  stesso premio in anni consecutivi e film differenti;
\item
  data di inizio di un prestito uguale alla fine di un altro.
\end{itemize}

\section{\texorpdfstring{14. ``Tutti'': divisione con doppio
\texttt{NOT\ EXISTS}}{14. ``Tutti'': divisione con doppio NOT EXISTS}}\label{14-tutti-divisione-con-doppio-not-exists}

``Trova gli \texttt{A} associati a \textbf{ogni} \texttt{B}
disponibile'' significa:

\begin{quote}
non esiste un B per il quale non esiste l\textquotesingle associazione
A-B.
\end{quote}

\begin{verbatim}
SELECT a.*
FROM A AS a
WHERE NOT EXISTS (
    SELECT 1
    FROM B AS b
    WHERE NOT EXISTS (
        SELECT 1
        FROM AB AS x
        WHERE x.aid = a.id
          AND x.bid = b.id
    )
);
\end{verbatim}

È il pattern per persone che hanno richiesto tutti i certificati e
artisti che hanno vinto in ogni categoria.

Alternativa con conteggi, valida se si gestiscono bene duplicati e
insieme vuoto:

\begin{verbatim}
HAVING COUNT(DISTINCT x.bid) = (SELECT COUNT(*) FROM B)
\end{verbatim}

\section{15. Conteggi distinti}\label{15-conteggi-distinti}

La parola ``diversi/e'' di solito richiede \texttt{DISTINCT} dentro
l\textquotesingle aggregato:

\begin{verbatim}
COUNT(DISTINCT certificate)
COUNT(DISTINCT venue)
\end{verbatim}

Non confondere \texttt{SELECT\ DISTINCT} (elimina righe duplicate del
risultato) con \texttt{COUNT(DISTINCT\ x)} (conta valori distinti di
\texttt{x}).

\section{16. Date}\label{16-date}

Per un anno usare un intervallo semiaperto, che funziona bene anche con
timestamp:

\begin{verbatim}
date_col >= DATE '2025-01-01'
AND date_col < DATE '2026-01-01'
\end{verbatim}

Per ``non scaduto'' la traccia del 14/01/2025 definisce esplicitamente
la validità come \texttt{expiry\_date\ \textgreater{}=\ CURRENT\_DATE};
non aggiungere altri requisiti se non richiesti.

\begin{center}\rule{0.5\linewidth}{0.5pt}\end{center}

Parte III - Algebra relazionale

\section{17. Operatori essenziali}\label{17-operatori-essenziali}

\begin{itemize}
\tightlist
\item
  selezione: \texttt{σ\_condizione(R)} - filtra tuple;
\item
  proiezione: \texttt{π\_attributi(R)} - sceglie attributi ed elimina
  duplicati;
\item
  ridenominazione: \texttt{ρ\_nuovoNome(R)} o
  \texttt{ρ\_\{nuovo←vecchio\}(R)};
\item
  prodotto: \texttt{R\ ×\ S};
\item
  theta/equi-join: \texttt{R\ ⋈\_condizione\ S};
\item
  unione \texttt{∪}, intersezione \texttt{∩}, differenza \texttt{−};
\item
  divisione \texttt{R(A,B)\ ÷\ S(B)} restituisce gli \texttt{A}
  associati a tutti i \texttt{B} di \texttt{S}.
\end{itemize}

Unione, intersezione e differenza richiedono compatibilità: stesso grado
e domini corrispondenti. Ridenominare gli attributi quando necessario.

\section{18. Pattern algebrici
d\textquotesingle esame}\label{18-pattern-algebrici-desame}

\subsection{Oggetti senza una proprietà
vietata}\label{oggetti-senza-una-proprietuxe0-vietata}

\begin{verbatim}
π_id(R) − π_rid(σ_condizione(S))
\end{verbatim}

Se la condizione richiede altre tabelle, prima costruire un join
ristretto.

\subsection{Almeno A e almeno B}\label{almeno-a-e-almeno-b}

\begin{verbatim}
π_id(σ_tipo='A'(X)) ∩ π_id(σ_tipo='B'(X))
\end{verbatim}

Non usare
\texttt{σ\_(tipo=\textquotesingle{}A\textquotesingle{}\ ∧\ tipo=\textquotesingle{}B\textquotesingle{})(X)}:
una singola tupla non può avere due valori diversi nello stesso
attributo.

\subsection{Tutti}\label{tutti}

\begin{verbatim}
π_{a,b}(AB) ÷ π_b(B)
\end{verbatim}

Gli attributi divisori devono avere lo stesso nome/dominio; usare
\texttt{ρ} se necessario.

\subsection{Massimo senza aggregati}\label{massimo-senza-aggregati}

Per trovare i valori non dominati:

\begin{verbatim}
T      = σ_condizione(R)
Piccoli = π_{t1.attributi}(
           σ_{t1.valore < t2.valore}(ρ_t1(T) × ρ_t2(T))
         )
Massimi = T − Piccoli
\end{verbatim}

Conserva tutti gli ex aequo.

\subsection{Ottimizzazione}\label{ottimizzazione}

\begin{enumerate}
\def\labelenumi{\arabic{enumi}.}
\tightlist
\item
  spingere le selezioni vicino alle relazioni di origine;
\item
  spingere le proiezioni, conservando gli attributi necessari ai join;
\item
  sostituire prodotti seguiti da selezione con join;
\item
  per un anti-join, sottrarre prima i soli identificatori.
\end{enumerate}

Esempio ``film senza premiere'':

\begin{verbatim}
Premiere = π_fid(σ_{is_premiere=true}(screening))
Senza    = π_fid(film) − Premiere
Risultato = π_{fid,title}(film ⋈ Senza)
\end{verbatim}

\begin{center}\rule{0.5\linewidth}{0.5pt}\end{center}

Parte IV - Progettazione ER e traduzione

\section{19. Cardinalità da una FK}\label{19-cardinalituxe0-da-una-fk}

Data \texttt{S.fk\ REFERENCES\ R(pk)}:

\begin{itemize}
\tightlist
\item
  ogni tupla di \texttt{S} è associata a \textbf{0..1 R} se la FK è
  nullable;
\item
  ogni tupla di \texttt{S} è associata a \textbf{1..1 R} se la FK è
  \texttt{NOT\ NULL};
\item
  una tupla di \texttt{R} può normalmente essere associata a
  \textbf{0..N S};
\item
  \texttt{UNIQUE(S.fk)} abbassa il massimo sul lato \texttt{R} a 1.
\end{itemize}

Una FK composta va trattata come un\textquotesingle unità concettuale.
Una FK ricorsiva crea un\textquotesingle associazione ricorsiva con
ruoli espliciti, per esempio \texttt{figlio} e \texttt{padre}.

\section{20. Da ER a relazionale}\label{20-da-er-a-relazionale}

\subsection{Entità}\label{entituxe0}

Ogni entità diventa una relazione; l\textquotesingle identificatore
diventa PK. Attributi opzionali diventano colonne nullable.

\subsection{Associazione 1:N}\label{associazione-1n}

Migrare la PK del lato 1 come FK nel lato N. La FK è nullable se la
partecipazione del lato N ha minimo 0.

\subsection{Associazione 1:1}\label{associazione-11}

Migrare la chiave in uno dei lati, preferibilmente in quello con
partecipazione totale, aggiungendo \texttt{UNIQUE} alla FK. Se
l\textquotesingle associazione ha attributi o si vuole mantenerla
autonoma, si può creare una relazione separata.

\subsection{Associazione N:M o n-aria}\label{associazione-nm-o-n-aria}

Creare una relazione con le FK verso tutte le entità partecipanti. Di
norma la PK è la combinazione delle FK. Gli attributi
dell\textquotesingle associazione restano qui.

Se la stessa coppia può ricomparire in date diverse,
l\textquotesingle associazione è storica: includere, per esempio,
\texttt{data\_inizio} nella PK.

\subsection{Attributo multivalore}\label{attributo-multivalore}

Creare una relazione separata, per esempio:

\begin{verbatim}
libro(bid, titolo)
genere(bid, genere)       PK(bid, genere), FK bid → libro.bid
\end{verbatim}

\subsection{Entità debole}\label{entituxe0-debole}

La PK contiene l\textquotesingle identificatore parziale più la PK
dell\textquotesingle entità proprietaria; quest\textquotesingle ultima è
anche FK.

\section{21. Generalizzazioni}\label{21-generalizzazioni}

\begin{itemize}
\tightlist
\item
  \textbf{Totale (T)}: ogni istanza del padre appartiene ad almeno un
  figlio.
\item
  \textbf{Parziale (P)}: possono esistere istanze del padre in nessun
  figlio.
\item
  \textbf{Esclusiva (E)}: un\textquotesingle istanza appartiene al
  massimo a un figlio.
\item
  \textbf{Sovrapposta (S)}: può appartenere a più figli.
\end{itemize}

Per padre \texttt{E(k,a)} e figli \texttt{E1(b)}, \texttt{E2(c)}:

\subsection{Strategia A - sola relazione
padre}\label{strategia-a---sola-relazione-padre}

\begin{verbatim}
E(k, a, b*, c*, tipo)
\end{verbatim}

\begin{itemize}
\tightlist
\item
  \texttt{tipo} è nullable se la gerarchia è parziale;
\item
  con gerarchia sovrapposta servono più indicatori o un valore per ogni
  combinazione;
\item
  servono vincoli extra-schema per coerenza tra tipo e attributi dei
  figli.
\end{itemize}

\subsection{Strategia B - sole relazioni
figlie}\label{strategia-b---sole-relazioni-figlie}

\begin{verbatim}
E1(k, a, b)
E2(k, a, c)
\end{verbatim}

È \textbf{impossibile per una gerarchia parziale}, perché perderebbe le
istanze del padre che non appartengono a figli. Con sovrapposizione
duplica i dati comuni.

\subsection{Strategia C - padre e
figli}\label{strategia-c---padre-e-figli}

\begin{verbatim}
E(k, a)
E1(k, b)     PK/FK k → E.k
E2(k, c)     PK/FK k → E.k
\end{verbatim}

È generale e riduce i nulli, ma richiede join. Totalità ed esclusività
spesso richiedono vincoli extra-schema.

\section{22. Da relazionale a ER:
minimizzazione}\label{22-da-relazionale-a-er-minimizzazione}

\begin{itemize}
\tightlist
\item
  Una tabella con propria PK e attributi descrittivi diventa normalmente
  entità.
\item
  Una tabella la cui PK è esattamente l\textquotesingle unione di due o
  più FK e che non ha attributi propri può essere minimizzata in
  un\textquotesingle associazione N:M.
\item
  Una tabella con PK propria e FK non identificanti resta entità
  collegata da associazioni.
\item
  Due FK dalla stessa tabella alla stessa entità sono due
  associazioni/ruoli distinti, non una sola.
\end{itemize}

\begin{center}\rule{0.5\linewidth}{0.5pt}\end{center}

Parte V - Normalizzazione

\section{23. Dipendenze funzionali e
chiusura}\label{23-dipendenze-funzionali-e-chiusura}

\texttt{X\ →\ Y} significa: tuple uguali su \texttt{X} devono essere
uguali anche su \texttt{Y}. Per cercare una chiave si calcola la
chiusura \texttt{X+}, aggiungendo ripetutamente gli attributi
determinati dalle DF. \texttt{X} è superchiave se \texttt{X+} contiene
tutti gli attributi; è chiave se è anche minimale.

Regole utili:

\begin{itemize}
\tightlist
\item
  riflessività: se \texttt{Y\ ⊆\ X}, allora \texttt{X\ →\ Y};
\item
  arricchimento: \texttt{X\ →\ Y} implica \texttt{XZ\ →\ YZ};
\item
  transitività: \texttt{X\ →\ Y} e \texttt{Y\ →\ Z} implicano
  \texttt{X\ →\ Z};
\item
  decomposizione: \texttt{X\ →\ YZ} implica \texttt{X\ →\ Y} e
  \texttt{X\ →\ Z};
\item
  unione: \texttt{X\ →\ Y} e \texttt{X\ →\ Z} implicano
  \texttt{X\ →\ YZ}.
\end{itemize}

\section{24. Forme normali}\label{24-forme-normali}

\subsection{1NF}\label{1nf}

Tutti gli attributi sono atomici; niente liste o attributi multivalore
nella stessa cella.

\subsection{2NF}\label{2nf}

Ogni attributo non primo dipende \textbf{completamente} da ogni chiave
candidata: niente dipendenze da una parte propria di una chiave
composta. Una relazione con tutte le chiavi semplici è automaticamente
in 2NF.

\subsection{3NF}\label{3nf}

Per ogni DF non banale \texttt{X\ →\ A}, almeno una condizione vale:

\begin{itemize}
\tightlist
\item
  \texttt{X} è una superchiave; oppure
\item
  \texttt{A} è primo.
\end{itemize}

\subsection{BCNF}\label{bcnf}

Per ogni DF non banale \texttt{X\ →\ A}, \texttt{X} deve essere una
superchiave. È più forte della 3NF.

Diagnosi tipica: con PK semplice \texttt{bid} e DF
\texttt{scaffale\ →\ ripiano}, la relazione è in 2NF ma non in 3NF/BCNF:
\texttt{scaffale} non è superchiave e \texttt{ripiano} non è primo.

\section{25. Decomposizione senza perdita e
conservazione}\label{25-decomposizione-senza-perdita-e-conservazione}

Una decomposizione è \textbf{senza perdita} se il join delle proiezioni
ricostruisce esattamente la relazione originaria, senza tuple spurie.
Per una decomposizione binaria \texttt{R\ →\ R1,R2}, è senza perdita se
l\textquotesingle intersezione \texttt{R1\ ∩\ R2} determina
funzionalmente tutti gli attributi di \texttt{R1} oppure tutti quelli di
\texttt{R2}.

È \textbf{conservativa rispetto alle dipendenze} se tutte le DF
originarie possono essere verificate sulle singole relazioni decomposte,
senza eseguire join.

La sintesi in 3NF può sempre ottenere perdita nulla e conservazione
delle DF. Una decomposizione BCNF può sacrificare la conservazione.

Procedura pratica d\textquotesingle esame:

\begin{enumerate}
\def\labelenumi{\arabic{enumi}.}
\tightlist
\item
  calcolare le chiavi;
\item
  trovare DF che violano la forma richiesta;
\item
  creare una relazione per ciascun determinante e i suoi dipendenti;
\item
  eliminare relazioni contenute in altre;
\item
  assicurarsi che almeno una relazione contenga una chiave originaria;
\item
  dichiarare PK e forma normale di ogni relazione.
\end{enumerate}

\begin{center}\rule{0.5\linewidth}{0.5pt}\end{center}

Parte VI - Sicurezza

\section{26. Politiche di accesso}\label{26-politiche-di-accesso}

\begin{itemize}
\tightlist
\item
  \textbf{Need-to-know / minimo privilegio}: accesso solo ai dati
  strettamente necessari. Ottima protezione, rischio di negare accessi
  innocui. Tipica di un sistema chiuso: tutto è vietato salvo
  autorizzazione esplicita.
\item
  \textbf{Maximized sharing / massima condivisione}: massimizza gli
  accessi, proteggendo solo ciò che è riservato. Più flessibile, meno
  prudente. Tipica di un sistema aperto: tutto è permesso salvo divieto
  esplicito.
\item
  \textbf{DAC, discrezionale}: proprietari/utenti possono concedere e
  revocare privilegi. Flessibile, ma non controlla il flusso dopo la
  lettura.
\item
  \textbf{MAC, mandatoria}: soggetti e oggetti hanno classi di
  sicurezza; le regole centrali controllano anche il flusso. Più rigida
  e adatta ad alta sicurezza.
\end{itemize}

System R è discrezionale, chiuso, con amministrazione decentralizzata
tramite ownership; supporta controllo per nome e per contenuto.

\section{\texorpdfstring{27. \texttt{GRANT}, grant option e
\texttt{REVOKE}}{27. GRANT, grant option e REVOKE}}\label{27-grant-grant-option-e-revoke}

\begin{verbatim}
GRANT SELECT, INSERT ON R TO utente WITH GRANT OPTION;
REVOKE SELECT ON R FROM utente CASCADE;
\end{verbatim}

\begin{itemize}
\tightlist
\item
  \texttt{WITH\ GRANT\ OPTION} permette al destinatario di delegare quel
  privilegio.
\item
  Si può delegare solo ciò che si possiede con grant option.
\item
  Ogni arco del grafo va etichettato con privilegio, tempo e presenza di
  GO.
\item
  Conviene disegnare un grafo separato per ogni privilegio.
\end{itemize}

Revoca ricorsiva: rimuovere l\textquotesingle arco revocato e tutti gli
archi che non avrebbero potuto essere creati senza di esso.
Un\textquotesingle altra sorgente salva il privilegio solo se è
indipendente e, per sostenere una delega già avvenuta, era disponibile
con grant option al tempo di quella delega. I timestamp contano.

\section{28. Cavallo di Troia}\label{28-cavallo-di-troia}

Un utente autorizzato esegue una procedura apparentemente lecita che
contiene codice nascosto. Sfruttando i privilegi
dell\textquotesingle esecutore, il codice legge un oggetto segreto e ne
copia il contenuto in un oggetto accessibile
all\textquotesingle attaccante. Le politiche discrezionali controllano
l\textquotesingle accesso iniziale ma non il successivo flusso
dell\textquotesingle informazione; una politica mandatoria di controllo
del flusso mira a impedirlo.

\begin{center}\rule{0.5\linewidth}{0.5pt}\end{center}

Parte VII - Transazioni e indici

\section{29. Transazioni e ACID}\label{29-transazioni-e-acid}

\begin{itemize}
\tightlist
\item
  \texttt{COMMIT}: termina con successo e rende permanenti gli effetti.
\item
  \texttt{ROLLBACK}: annulla gli effetti non confermati.
\item
  In \texttt{AUTOCOMMIT} ogni istruzione SQL è una transazione.
\end{itemize}

ACID:

\begin{itemize}
\tightlist
\item
  \textbf{Atomicity}: tutto o niente.
\item
  \textbf{Consistency}: da uno stato che soddisfa i vincoli a un altro
  stato che li soddisfa.
\item
  \textbf{Isolation}: esecuzione concorrente equivalente a un ordine
  seriale.
\item
  \textbf{Durability}: dopo il commit gli effetti persistono anche dopo
  guasti.
\end{itemize}

Anomalie:

\begin{itemize}
\tightlist
\item
  \textbf{lost update}: un aggiornamento sovrascrive quello concorrente;
\item
  \textbf{dirty read}: T2 legge un dato scritto da T1 prima del commit;
  T1 poi fa rollback e T2 ha usato un valore mai esistito stabilmente;
\item
  \textbf{incorrect summary}: un aggregato legge una parte dei dati
  prima e una parte dopo un aggiornamento concorrente.
\end{itemize}

Il dirty read si previene non rendendo visibili modifiche non
confermate, con livello almeno \texttt{READ\ COMMITTED} e/o lock
appropriati.

\section{30. Strutture fisiche e
indici}\label{30-strutture-fisiche-e-indici}

\begin{itemize}
\tightlist
\item
  Strutture primarie: sequenziale, hash (accesso calcolato), albero.
\item
  Hash: molto efficiente per uguaglianza, non per intervalli; gestisce
  collisioni. Nel lessico delle dispense è una struttura primaria.
\item
  Un indice primario determina la disposizione dei record ed è
  \textbf{sparso}: una voce per blocco dati.
\item
  Un indice secondario non determina la disposizione ed è sempre
  \textbf{denso}: deve rappresentare ogni valore/record indicizzato.
\end{itemize}

Formule:

\begin{verbatim}
bfr = floor(B / R)
blocchi_dati = ceil(numero_record / bfr)
voci_indice_primario = blocchi_dati
\end{verbatim}

Esempio d\textquotesingle esame: 20.000 record, \texttt{bfr\ =\ 10} →
2.000 blocchi → 2.000 voci nell\textquotesingle indice primario.

\begin{center}\rule{0.5\linewidth}{0.5pt}\end{center}

Parte VIII - Ripasso minimo degli argomenti non comparsi

\section{31. DBMS, livelli e
indipendenza}\label{31-dbms-livelli-e-indipendenza}

Un DBMS gestisce dati persistenti, condivisi e consistenti, offrendo
controllo di concorrenza, sicurezza, recupero e linguaggi di
interrogazione. I livelli esterno, concettuale e interno separano viste
utente, schema logico e memorizzazione fisica.
L\textquotesingle indipendenza logica limita l\textquotesingle impatto
dei cambiamenti concettuali sulle viste; quella fisica limita
l\textquotesingle impatto dei cambiamenti di memorizzazione sullo schema
logico.

\section{32. Trigger e viste}\label{32-trigger-e-viste}

Un trigger segue il paradigma evento-condizione-azione e reagisce a
\texttt{INSERT}, \texttt{UPDATE} o \texttt{DELETE},
\texttt{BEFORE}/\texttt{AFTER}, per riga o per istruzione. È utile per
vincoli non esprimibili nel DDL e dati derivati, ma va evitato quando un
vincolo dichiarativo basta.

Una vista è una query con nome. Una vista semplice su una sola tabella
può essere aggiornabile; join, aggregazioni e raggruppamenti in genere
impediscono l\textquotesingle aggiornamento diretto.
\texttt{WITH\ CHECK\ OPTION} impedisce modifiche che rendano la riga
invisibile nella vista.

\section{33. NoSQL in una frase}\label{33-nosql-in-una-frase}

I sistemi NoSQL rinunciano a parte dell\textquotesingle uniformità del
modello relazionale per flessibilità, distribuzione o specifici pattern
di accesso. Nei document store come MongoDB i dati sono documenti
JSON/BSON annidati; la scelta tra embedding e riferimenti dipende da
cardinalità, aggiornamenti e letture congiunte.

\begin{center}\rule{0.5\linewidth}{0.5pt}\end{center}

Checklist finale da memorizzare

\begin{itemize}
\tightlist
\item
  PK = \texttt{UNIQUE\ +\ NOT\ NULL}; ogni colonna della PK composta è
  non nulla.
\item
  FK senza \texttt{NOT\ NULL} può essere nulla; FK senza
  \texttt{CASCADE} blocca la modifica.
\item
  \texttt{LEFT\ JOIN} + \texttt{WHERE\ destra.id\ IS\ NULL} = ``senza''.
\item
  Filtro sul lato destro nell\textquotesingle{}\texttt{ON} per non
  perdere gli zeri.
\item
  \texttt{COUNT(colonna\_destra)}, non \texttt{COUNT(*)}, dopo un
  \texttt{LEFT\ JOIN}.
\item
  ``diversi'' → spesso \texttt{COUNT(DISTINCT\ ...)}.
\item
  ``tutti'' → divisione o doppio \texttt{NOT\ EXISTS}.
\item
  ``massimo'' → confrontare con \texttt{MAX}, conservando gli ex aequo.
\item
  coppie → self join; \texttt{\textless{}} elimina simmetrie,
  \texttt{+\ 1} esprime anni consecutivi.
\item
  algebra: ``almeno A e almeno B'' → intersezione, non
  \texttt{A\ AND\ B} sulla stessa tupla.
\item
  2NF elimina dipendenze parziali; 3NF ammette RHS primo; BCNF pretende
  determinante superchiave.
\item
  revoca: separare i privilegi e rispettare provenienza, grant option e
  tempo.
\item
  indice primario sparso, secondario denso; \texttt{bfr\ =\ floor(B/R)}.
\item
  \texttt{COMMIT} successo; coerenza = rispetto dei vincoli; dirty read
  = lettura di scritture non confermate.
\end{itemize}

\end{document}
